\documentclass[preview]{standalone}

\usepackage{amsmath}
\usepackage{amssymb}
\usepackage{bettelini}
\usepackage{stellar}
\usepackage{definitions}

\begin{document}

\id{measuretheory-lp-spaces}
\genpage

\section{Lebesgue Spaces}

\begin{snippetdefinition}{l-p-space-definition}{The space \(L^p\)}
    Let \((X,\mathcal M,\mu)\) be a measure space and
    \(1\leq p<+\infty\). Let
    \[
        \mathcal L^p(\mu)
        =
        \left\{
            f\colon X\fromto\complexnumbers\text{ measurable}
            \suchthat \int_X|f|^p\,d\mu<+\infty
        \right\}.
    \]
    After identifying functions that agree almost everywhere, define
    \[
        L^p(\mu)=\mathcal L^p(\mu)/{\sim},
        \qquad
        \|f\|_p=
        \left(\int_X|f|^p\,d\mu\right)^{1/p}.
    \]
    For \(p=+\infty\), define
    \[
        \|f\|_\infty=\operatorname*{ess\,sup}_X|f|.
    \]
\end{snippetdefinition}

\begin{snippettheorem}{l-p-space-banach-theorem}{Riesz--Fischer theorem}
    For every measure space and every \(1\leq p\leq+\infty\),
    \((L^p(\mu),\|\cdot\|_p)\) is a Banach space.
\end{snippettheorem}

\begin{snippetexample}{discrete-l-p-spaces-example}{Discrete \(L^p\) spaces}
    If \(X=\{1,\ldots,n\}\) has counting measure, then
    \(L^p(\mu)=\realnumbers^n\) or \(\complexnumbers^n\), with
    \[
        \|x\|_p=\left(\sum_{i=1}^n|x_i|^p\right)^{1/p}.
    \]
    For \(p=2\), this is the Euclidean norm. On
    \(X=\naturalnumbers\) with counting measure, one obtains the sequence
    spaces \(\ell^p\); in particular, \(\ell^1\) consists of absolutely
    summable sequences and \(\ell^2\) of square-summable sequences.
\end{snippetexample}

\section{Duality and the Integral Functional}

\begin{snippetproposition}{integral-continuous-functional-on-l-one}{The integral is continuous on \(L^1\)}
    The map
    \[
        I\colon L^1(\mu)\fromto\complexnumbers,
        \qquad
        I(f)=\int_Xf\,d\mu,
    \]
    is a continuous linear functional, and \(\|I\|\leq1\).
\end{snippetproposition}

\begin{snippetproof}{integral-continuous-functional-on-l-one-proof}{integral-continuous-functional-on-l-one}{The integral is continuous on \(L^1\)}
    Linearity is inherited from the integral, and
    \[
        |I(f)|=\left|\int_Xf\,d\mu\right|
        \leq\int_X|f|\,d\mu=\|f\|_1.
    \]
    Hence \(I\) is continuous. Equivalently, if \(f_n\to f\) in \(L^1\),
    then \(I(f_n)\to I(f)\).
\end{snippetproof}

\begin{snippet}{l-p-duality-overview}
    In infinite-dimensional vector spaces, an algebraic linear functional
    need not be continuous, so the algebraic dual can be strictly larger than
    the continuous dual. For \(1<p<+\infty\), under standard hypotheses such
    as \(\sigma\)-finiteness, the continuous dual of \(L^p(\mu)\) is
    identified with \(L^q(\mu)\), where
    \[
        \frac1p+\frac1q=1.
    \]
    In particular, \(L^2(\mu)\) is identified with its continuous dual.
\end{snippet}

\section{Convergence in L1}

\begin{snippet}{continuous-representatives-in-l-one}
    An \(L^1(\realnumbers)\) element is an equivalence class. Some classes
    contain a continuous representative and others do not. If a continuous
    representative exists, it is unique: two distinct continuous functions
    differ on a nonempty open set, which has positive Lebesgue measure. The
    Dirichlet function is not a counterexample because it represents the zero
    class. Fat Cantor sets provide classes with no Riemann-integrable
    representative.
\end{snippet}

\begin{snippettheorem}{riesz-subsequence-principle}{Riesz subsequence principle}
    If \(f_n\to f\) in \(L^1(\mu)\), then there is a subsequence
    \(f_{n_k}\to f\) almost everywhere. Moreover, if any subsequence
    \(f_{n_k}\) converges almost everywhere to \(g\), then \(f=g\) almost
    everywhere.
\end{snippettheorem}

\begin{snippetproof}{riesz-subsequence-principle-proof}{riesz-subsequence-principle}{Riesz subsequence principle}
    Choose a subsequence such that
    \[
        \|f_{n_k}-f\|_1<2^{-k}.
    \]
    Then
    \[
        \sum_{k=1}^{\infty}\int_X|f_{n_k}-f|\,d\mu<+\infty.
    \]
    The theorem on integration of absolutely summable series implies that
    \(\sum_k|f_{n_k}-f|<+\infty\) almost everywhere. Hence
    \(|f_{n_k}-f|\to0\) almost everywhere.

    Now suppose that a subsequence \(f_{n_k}\to g\) almost everywhere. It
    still converges to \(f\) in \(L^1\), so the first part gives a further
    subsequence converging almost everywhere to \(f\). The same further
    subsequence also converges almost everywhere to \(g\), and uniqueness of
    pointwise limits gives \(f=g\) almost everywhere.
\end{snippetproof}

\begin{snippetproposition}{essential-uniform-convergence-implies-l-one}{Essentially uniform convergence on a finite measure space}
    Let \(\mu(X)<+\infty\). If
    \[
        \|f_n-f\|_\infty
        =\operatorname*{ess\,sup}_X|f_n-f|
        \longrightarrow0,
    \]
    then \(f_n\to f\) in \(L^1(\mu)\).
\end{snippetproposition}

\begin{snippetproof}{essential-uniform-convergence-implies-l-one-proof}{essential-uniform-convergence-implies-l-one}{Essentially uniform convergence on a finite measure space}
    Directly,
    \[
        \|f_n-f\|_1
        =\int_X|f_n-f|\,d\mu
        \leq\mu(X)\|f_n-f\|_\infty\longrightarrow0.
    \]
\end{snippetproof}

\section{Density and Size}

\begin{snippettheorem}{integrable-simple-functions-dense-in-l-one}{Integrable simple functions are dense in \(L^1\)}
    Let \(S\) be the set of measurable simple functions \(s\) such that
    \[
        \mu(\{x\in X\suchthat s(x)\neq0\})<+\infty.
    \]
    Then \(S\) is dense in \(L^1(\mu)\).
\end{snippettheorem}

\begin{snippetproof}{integrable-simple-functions-dense-in-l-one-proof}{integrable-simple-functions-dense-in-l-one}{Integrable simple functions are dense in \(L^1\)}
    If \(f\geq0\) is integrable, choose integrable simple functions
    \(s_n\uparrow f\). Monotone convergence gives
    \[
        \|f-s_n\|_1
        =\int_X(f-s_n)\,d\mu
        \longrightarrow0.
    \]
    A real or complex function is reduced to this case by decomposing its
    real and imaginary parts into positive and negative parts.
\end{snippetproof}

\begin{snippet}{l-one-separability-and-dimension}
    In many standard measure spaces, for example when the
    \(\sigma\)-algebra is generated modulo null sets by a countable family,
    \(L^1(\mu)\) is separable. Separability is not automatic for an arbitrary
    measure space.

    If there are infinitely many pairwise disjoint measurable sets of
    positive finite measure, their indicator functions are linearly
    independent, so \(L^1(\mu)\) is infinite-dimensional. The sequence space
    \(\ell^1\) is a basic infinite-dimensional example.
\end{snippet}

\begin{snippetproposition}{nontrivial-banach-space-uncountable}{A nontrivial Banach space is uncountable}
    A nontrivial Banach space cannot have countably infinite cardinality.
\end{snippetproposition}

\begin{snippetproof}{nontrivial-banach-space-uncountable-proof}{nontrivial-banach-space-uncountable}{A nontrivial Banach space is uncountable}
    Suppose that \(X=\{x_1,x_2,\ldots\}\). Every singleton is closed and has
    empty interior: if \(B(x_n,r)\subseteq\{x_n\}\), choose a nonzero
    \(v\in X\) and sufficiently small \(t>0\); then
    \(x_n+tv\in B(x_n,r)\) but \(x_n+tv\neq x_n\). Thus every singleton is
    nowhere dense and
    \[
        X=\bigcup_{n=1}^{\infty}\{x_n\}
    \]
    is meagre in itself, contradicting the Baire category theorem for the
    complete metric space \(X\).
\end{snippetproof}

\begin{snippet}{l-one-tail-behavior}
    For \(f\in L^1(\realnumbers)\), most of the integral can be confined to a
    finite-measure region. Nevertheless, a representative need not satisfy
    \(f(x)\to0\) as \(|x|\to+\infty\): one may alter it at every integer
    without changing its class, and even continuous integrable functions can
    have narrow peaks of fixed height escaping to infinity. One can still
    choose points \(x_n\to+\infty\) along which a suitable representative
    tends to zero.
\end{snippet}

\end{document}
